% =====================================================================================================
% PART 2 (frames 7-9): background -- Hastad's theorem, gadget reductions, the colour-reference problem.
% Sources in the paper: Section 1.2 (first paragraph), Section 2 (Theorem 2.1, gadgets, the reduction,
% Figure 2).
% =====================================================================================================

% ---- Frame 7: Starting point: Håstad's theorem
\begin{frame}{Starting point: H\aa stad's theorem}
\vspace{2mm}
\begin{columns}[c,onlytextwidth]
\begin{column}{0.55\textwidth}
Linear equations over $\Z_3$ with three variables:\\[2.6mm]
\hspace*{12mm}{\Large$x_1+x_2+x_3=t$}\\[2.6mm]
$t\in\{0,1,2\}$ is the \hl{type} of the equation
\end{column}
\begin{column}{0.42\textwidth}
\raggedleft
\begin{tikzpicture}
\node[draw=black!30,fill=black!3,rounded corners=3pt,inner xsep=7pt,inner ysep=4.5pt] (ex) {%
  \renewcommand{\arraystretch}{1.12}%
  $\begin{array}{@{}r@{\;}c@{\;}l@{\qquad}l@{}}
   {\color{acc}u_1}+w_1+w_2 &=& 0 & \text{\color{soft}\small type $0$}\\
   {\color{acc}u_1}+w_2+w_3 &=& 1 & \text{\color{soft}\small type $1$}\\
   {\color{acc}u_2}+w_3+w_4 &=& 2 & \text{\color{soft}\small type $2$}
  \end{array}$};
\node[anchor=south west,font=\footnotesize,text=soft,inner sep=1.5pt] at (ex.north west) {example};
\end{tikzpicture}
\end{column}
\end{columns}
\vspace{3mm}
\begin{block}{Theorem 2.1 \src{H\aa s01}}
For every $\eps>0$ it is NP-hard to distinguish\\[1mm]
\begin{tabular}{@{\hspace{3mm}}l@{\quad}l@{}}
\hl{YES} & some assignment satisfies $\ge 1-\eps$ of the equations\\[0.6mm]
\bad{\textbf{NO}} & every assignment satisfies $\le \tfrac13+\eps$ of the equations
\end{tabular}%
\hfill\raisebox{-2.55mm}{\color{soft}\small random assignment: $\tfrac13$}
\end{block}
\vspace{3mm}
{\small
\begin{tabular}{@{}l@{\quad}l@{}}
We use two features & {\tiny\color{acc}\raisebox{0.25ex}{$\blacksquare$}}\ \ \hl{role separation}: one variable from
  ${\color{acc}U}$, two from $W$ \ ($U$, $W$ disjoint)\\[0.3mm]
of these instances: &{\tiny\color{acc}\raisebox{0.25ex}{$\blacksquare$}}\ \ \hl{balance}: each type has total weight $\tfrac13$
\end{tabular}}
\end{frame}

% ---- Frame 8: Gadget reductions
\begin{frame}{Gadget reductions}
\vspace{0.5mm}
\begin{columns}[c,onlytextwidth]
\begin{column}{0.64\textwidth}
% Frame 8: schematic of the reduction, adapted from the paper's Figure 2 (paper/figures/reduction.tex).
% Same instance as in the paper: U = {u_1,u_2}, W = {w_1,...,w_4}, three equations
%   u_1+w_1+w_2=0,  u_1+w_2+w_3=1,  u_2+w_3+w_4=2   (one of each type),
% each replaced by a copy of the gadget I_t (box; new auxiliary vertices inside, schematic) whose three
% primaries are identified with the three variable vertices of the equation (grey legs).
% Simplified for the slide: larger labels; the row of global vertices, the labels p_1,p_2,p_3 and the
% "decoded by a / b" labels of the paper's figure are dropped.
\begin{tikzpicture}[x=1cm,y=1cm,
  prim/.style={circle,draw,fill=black!15,inner sep=0.5pt,minimum size=17pt,font=\small,line width=0.5pt},
  aux/.style={circle,draw,fill=white,inner sep=0pt,minimum size=5.6pt,line width=0.45pt},
  gad/.style={draw=acc!55,fill=acc!5,rounded corners=4pt,line width=0.6pt,
              minimum width=4.9cm,minimum height=0.82cm,inner sep=0pt},
  leg/.style={line width=0.6pt,black!55},
  gname/.style={font=\small,inner sep=1pt,text=acc!90!black},
  eqn/.style={font=\small,inner sep=1pt,anchor=west},
  clab/.style={font=\normalsize,text=black!70,inner sep=1pt},
  note/.style={font=\footnotesize,text=soft,inner sep=1.5pt},
  ptr/.style={-{Stealth[length=3.5pt,width=3pt]},line width=0.4pt,soft}]
\def\xu{0}\def\xw{8.4}\def\bx{4.2}% U column, W column, centre of the gadget copies
% ---------------------------------------------------------------- variables (one vertex each)
\node[prim] (u1) at (\xu,2.30) {$u_1$};
\node[prim] (u2) at (\xu,0.575) {$u_2$};
\node[prim] (w1) at (\xw,3.45) {$w_1$};
\node[prim] (w2) at (\xw,2.30) {$w_2$};
\node[prim] (w3) at (\xw,1.15) {$w_3$};
\node[prim] (w4) at (\xw,0.00) {$w_4$};
% ---------------------------------------------------------------- gadget copies
\foreach \n/\y/\t/\eq in {0/2.875/0/{u_1+w_1+w_2=0}, 1/1.725/1/{u_1+w_2+w_3=1}, 2/0.575/2/{u_2+w_3+w_4=2}}{
  \node[gad] (B\n) at (\bx,\y) {};
  \node[gname] at ([xshift=0.40cm]B\n.west) {$I_{\t}$};
  \draw[acc!30,line width=0.4pt]
    ([xshift=0.76cm,yshift=0.23cm]B\n.west) -- ([xshift=0.76cm,yshift=-0.23cm]B\n.west);
  \node[eqn] at ([xshift=0.86cm]B\n.west) {$\eq$};
  \node[aux] (x\n a) at ([xshift=-0.96cm,yshift=0.13cm]B\n.east) {};
  \node[aux] (x\n b) at ([xshift=-0.62cm,yshift=-0.13cm]B\n.east) {};
  \node[aux] (x\n c) at ([xshift=-0.28cm,yshift=0.13cm]B\n.east) {};
  \draw[black!55,line width=0.4pt] (x\n a) -- (x\n b) -- (x\n c) -- (x\n a);}
% ---------------------------------------------------------------- legs: first variable in U, the other two in W
\draw[leg] (B0.west) -- (u1);
\draw[leg] (B1.west) -- (u1);
\draw[leg] (B2.west) -- (u2);
\draw[leg] ([yshift=0.18cm]B0.east) -- (w1);
\draw[leg] ([yshift=-0.18cm]B0.east) -- (w2);
\draw[leg] ([yshift=0.18cm]B1.east) -- (w2);
\draw[leg] ([yshift=-0.18cm]B1.east) -- (w3);
\draw[leg] ([yshift=0.18cm]B2.east) -- (w3);
\draw[leg] ([yshift=-0.18cm]B2.east) -- (w4);
% ---------------------------------------------------------------- column labels
\node[clab,anchor=base] at (\xu,3.98) {$U$};
\node[clab,anchor=base] at (\xw,3.98) {$W$};
% ---------------------------------------------------------------- annotations
\node[note,anchor=base] (lg) at (\bx-1.25,3.98) {gadget $I_t$ for type $t$};
\draw[ptr] (lg.south) -- ([xshift=-1.25cm,yshift=0.03cm]B0.north);
\node[note,anchor=base] (la) at ([xshift=-0.62cm]B0.east |- 0,3.98) {auxiliary vertices};
\draw[ptr] (la.south) -- ([yshift=0.07cm]x0b.north |- x0a.north);
\end{tikzpicture}

\end{column}
\begin{column}{0.34\textwidth}
\begin{itemize}\setlength{\itemsep}{3.2mm}
\item Variable $\to$ \hl{one vertex};\\ its colour = its value
\item Equation $\to$ a \hl{gadget}:\\ a small graph with new\\ auxiliary vertices
\item Equation satisfied $\Rightarrow$\\ the gadget can cut more
\end{itemize}
\end{column}
\end{columns}
\vspace{1mm}
\begin{block}{}
\hspace*{\fill}%
\begin{tabular}{@{}r@{\;\;$\longrightarrow$\;\;}l@{}}
almost satisfiable system & \hl{large cut}\\[0.5mm]
far from satisfiable & \bad{\textbf{every colouring cuts less}}
\end{tabular}%
\hspace*{\fill}\llap{\raisebox{-2.3mm}{\src{TSSW00}}}
\end{block}
\end{frame}

% ---- Frame 9: The obstacle: colours have no names
\begin{frame}{The obstacle: colours have no names}
\vspace{3mm}
\begin{columns}[T,onlytextwidth]
% ---------------- zone (i): renaming the colours
\begin{column}{0.315\textwidth}
\textbf{Rename the colours:}\par\vspace{3.4mm}
\centerline{% Frame 9, zone (i): the same 3-coloured graph twice (a triangle with a pendant edge); in the right copy
% the grey and the blue colour are exchanged everywhere.  The same three edges are cut and the same edge
% (dashed, warn) is not cut.  Colours as fills: white, grey (black!35), accent blue.
\begin{tikzpicture}[x=1cm,y=1cm,
  v/.style={circle,draw,line width=0.6pt,inner sep=0pt,minimum size=11pt},
  ca/.style={fill=white}, cb/.style={fill=black!35}, cc/.style={fill=acc},
  e/.style={line width=0.8pt},
  unc/.style={warn,line width=1.8pt,dash pattern=on 3pt off 1.6pt},
  tip/.tip={Stealth[length=4pt,width=3.6pt]}]
% ---- left copy: triangle (white, grey, blue) and a pendant vertex with the colour of its neighbour
\node[v,ca] (a) at (0.20,1.30) {};
\node[v,cb] (b) at (0,0) {};
\node[v,cc] (c) at (1.10,0) {};
\node[v,cc] (d) at (1.10,1.30) {};
\draw[e] (a) -- (b) -- (c) -- (a);
\draw[unc] (c) -- (d);
% ---- right copy: grey and blue exchanged
\begin{scope}[shift={(2.85,0)}]
  \node[v,ca] (a2) at (0.20,1.30) {};
  \node[v,cc] (b2) at (0,0) {};
  \node[v,cb] (c2) at (1.10,0) {};
  \node[v,cb] (d2) at (1.10,1.30) {};
  \draw[e] (a2) -- (b2) -- (c2) -- (a2);
  \draw[unc] (c2) -- (d2);
\end{scope}
% ---- the renaming: exchange grey and blue
\draw[tip-tip,line width=0.8pt,acc] (1.55,0.42) -- (2.40,0.42);
\node[v,cb,minimum size=7.5pt,line width=0.45pt] (s1) at (1.66,0.90) {};
\node[v,cc,minimum size=7.5pt,line width=0.45pt] (s2) at (2.29,0.90) {};
\draw[{Stealth[length=2.6pt,width=2.4pt]}-{Stealth[length=2.6pt,width=2.4pt]},line width=0.45pt]
  ([xshift=1.2pt]s1.east) -- ([xshift=-1.2pt]s2.west);
\end{tikzpicture}
}\par\vspace{2.6mm}
\centerline{exactly the \hl{same cut}}
\end{column}
\begin{column}{0.4pt}\color{acc!25}\rule{0.4pt}{37mm}\end{column}
% ---------------- zone (ii): type 0
\begin{column}{0.285\textwidth}
\textbf{Type 0 is fine:}\par\vspace{2.4mm}
$x_1+x_2+x_3=0$\quad$\Longleftrightarrow$\par\vspace{3mm}
\hspace*{2mm}% Frame 9, zone (ii): the two kinds of colour patterns with x_1+x_2+x_3 = 0 over Z_3:
% all three values equal, or all three distinct.  Colours as fills: white, grey (black!35), accent blue.
\begin{tikzpicture}[x=1cm,y=1cm,
  v/.style={circle,draw,line width=0.6pt,inner sep=0pt,minimum size=10pt},
  ca/.style={fill=white}, cb/.style={fill=black!35}, cc/.style={fill=acc},
  lab/.style={font=\small,anchor=west,inner sep=0pt}]
\foreach \i/\c in {0/cb,1/cb,2/cb}{\node[v,\c] at (0.46*\i,0.62) {};}
\node[lab] at (1.38,0.62) {all equal, or};
\foreach \i/\c in {0/ca,1/cb,2/cc}{\node[v,\c] at (0.46*\i,0) {};}
\node[lab] at (1.38,0) {all distinct};
\end{tikzpicture}
\par\vspace{2.6mm}
{\small unchanged when the colours are renamed}
\end{column}
\begin{column}{0.4pt}\color{acc!25}\rule{0.4pt}{37mm}\end{column}
% ---------------- zone (iii): types 1 and 2
\begin{column}{0.335\textwidth}
\textbf{Types 1 and 2 are not:}\par\vspace{2.4mm}
which colour is ``$1$''?\par\vspace{1.8mm}
$\Rightarrow$ global \hl{reference vertices}\par\vspace{3mm}
\hspace*{4mm}% Frame 9, zone (iii): the global reference vertices g_0,g_1,g_2 (squares).
% Top row: an honest colouring gives them three different colours (they play the roles of 0,1,2).
% Bottom row (warn): a cheating colouring gives all three the same colour.
% Colours as fills: white, grey (black!35), accent blue.
\begin{tikzpicture}[x=1cm,y=1cm,
  g/.style={rectangle,draw,line width=0.6pt,inner sep=0pt,minimum size=10pt},
  ca/.style={fill=white}, cb/.style={fill=black!35}, cc/.style={fill=acc},
  rl/.style={font=\small,anchor=east,inner sep=0pt},
  gl/.style={font=\small,inner sep=0pt,anchor=base}]
\foreach \i/\x in {0/0,1/0.78,2/1.56}{\node[gl] at (\x,0.98) {$g_{\i}$};}
\node[rl] at (-0.42,0.58) {honest};
\foreach \x/\c in {0/ca,0.78/cb,1.56/cc}{\node[g,\c] at (\x,0.58) {};}
\node[rl,text=warn] at (-0.42,0) {cheating};
\foreach \x in {0,0.78,1.56}{\node[g,cb,draw=warn,line width=1.1pt] at (\x,0) {};}
\end{tikzpicture}

\end{column}
\end{columns}
\vspace{4.5mm}
\begin{alertblock}{The colour-reference problem}
Guarding against cheating makes the gadgets for $t\neq0$ \bad{\textbf{large and inefficient}}:
for $k=3$, the gadgets of \src{NRT25} give only $55/57\approx0.965$, weaker than
$16/17\approx0.941$.
\end{alertblock}
\end{frame}
